\chapter*{Conclusion Générale et Perspectives}
\addcontentsline{toc}{chapter}{Conclusion Générale et Perspectives}
\markboth{Conclusion Générale et Perspectives}{}

\section*{Conclusion Générale}

Ce projet de fin d'études, réalisé au sein de l'entreprise \textbf{TAMTAM International}, m'a permis de concevoir et de développer une application de gestion des webinaires destinée aux fiduciaires belges. Tout au long de ce stage, j'ai pu mettre en pratique les connaissances théoriques acquises durant ma formation en Master Ingénierie des Systèmes d'Information, tout en me confrontant aux réalités et aux exigences du monde professionnel.

Sur le plan technique, ce projet m'a permis de maîtriser un ensemble de technologies modernes et complémentaires : \textbf{Meteor.js} et \textbf{React.js} côté frontend, \textbf{Symfony} côté backend, ainsi que les bases de données \textbf{MongoDB} et \textbf{MySQL}. La mise en œuvre de la méthodologie \textbf{Agile Scrum} a favorisé une organisation rigoureuse du travail, une livraison itérative des fonctionnalités et une communication efficace au sein de l'équipe.

Les résultats obtenus sont satisfaisants : la solution développée permet de gérer efficacement les webinaires, d'automatiser plusieurs opérations clés — telles que l'envoi de certificats de participation et la gestion des sondages — et d'améliorer sensiblement l'expérience des utilisateurs. Les fonctionnalités livrées répondent aux besoins exprimés par les parties prenantes et s'intègrent harmonieusement dans l'écosystème existant de la plateforme TAMTAM.

Sur le plan personnel, cette expérience a été particulièrement enrichissante. Elle m'a permis de développer mon autonomie, mon sens des responsabilités et ma capacité à travailler en équipe dans un contexte professionnel international. Elle m'a également conforté dans mon choix de carrière en ingénierie des systèmes d'information.

\section*{Perspectives}

Bien que les objectifs initiaux aient été atteints, plusieurs pistes d'amélioration et d'évolution méritent d'être envisagées pour les versions futures de la plateforme :

\begin{itemize}
    \item \textbf{Amélioration des performances} : optimisation des requêtes et mise en cache pour gérer un plus grand nombre d'utilisateurs simultanés lors des webinaires en direct.

    \item \textbf{Extension des fonctionnalités mobiles} : développement d'une application mobile native (iOS/Android) pour offrir une expérience utilisateur optimale sur smartphones et tablettes.

    \item \textbf{Intelligence artificielle} : intégration de mécanismes de recommandation basés sur le comportement des utilisateurs pour suggérer des webinaires pertinents et personnaliser le contenu.

    \item \textbf{Accessibilité} : mise en conformité avec les standards d'accessibilité WCAG 2.1 afin de rendre la plateforme utilisable par des personnes en situation de handicap.

    \item \textbf{Internationalisation} : extension du support multilingue pour couvrir de nouveaux marchés au-delà de la Belgique francophone.

    \item \textbf{Analytique avancée} : mise en place de tableaux de bord analytiques permettant aux organisateurs de suivre en temps réel les indicateurs clés de performance de leurs webinaires.
\end{itemize}

Ces perspectives ouvrent des horizons prometteurs pour la plateforme TAMTAM et témoignent du potentiel d'évolution d'une solution déjà robuste et fonctionnelle.
